% chapters/04/checklist.tex
\bagian{Checklist Pencapaian Kompetensi}

Bab ini menuntaskan \SUBCPMK{1.1}. Centang setiap butir berikut hanya apabila
Anda benar-benar telah melakukannya atau dapat menjelaskannya.

\begin{checklistbox}
\begin{checklist}
  \item {\bfseries (Indikator a)} Saya dapat menyebutkan ketiga tahap siklus
        permainan secara berurutan beserta pekerjaan masing-masing.
  \item Saya dapat menjelaskan mengapa masukan dibaca lebih dahulu dan
        penayangan dikerjakan paling akhir.
  \item {\bfseries (Indikator b)} Saya dapat menjelaskan dua pekerjaan
        \code{clock.tick()} sekaligus.
  \item Saya dapat menjelaskan mengapa \code{clock.tick()} diletakkan pada
        akhir putaran siklus.
  \item Saya dapat membatasi laju bingkai program saya sendiri dan
        menampilkan laju terukurnya pada judul jendela.
  \item Saya dapat mengukur pengaruh batas laju terhadap pemakaian prosesor
        dan menjelaskan hasilnya.
  \item Saya dapat menjelaskan mengapa siklus tanpa batas laju memboroskan
        prosesor.
  \item Saya dapat menjelaskan mengapa menaikkan batas laju tidak selalu
        menaikkan laju yang tercapai.
  \item {\bfseries (Indikator c)} Saya dapat membedakan langkah tetap dan
        langkah berubah, serta menentukan golongan mana yang dipakai siklus
        pada bab ini.
  \item Saya dapat menjelaskan akibat panjang langkah yang terlalu besar
        terhadap deteksi tabrakan.
  \item Program siklus saya berjalan terus sampai jendela ditutup dan berakhir
        dengan rapi, bukan dihentikan paksa.
  \item Saya dapat menuliskan tiga gejala yang muncul pada siklus yang salah
        beserta penyebabnya.
  \item Catatan praktikum saya memuat hasil pengukuran batas laju pada komputer
        saya sendiri.
\end{checklist}
\end{checklistbox}
